% PAGES: 24-30, 31-38

\section{Matrix rank, nullspace, and range of a matrix}

\begin{definition}
Let $w_1, w_2, \ldots, w_p$ be vectors of the same size. The vectors $w_1,
w_2, \ldots, w_p$ are linearly independent if and only if the only linear
combination of these vectors that is equal to 0 has all coefficients equal
to 0, i.e.,
\[
\text{if} \quad \sum_{i=1}^p c_iw_i = 0, \quad \text{with}\ c_i\in\R,\
i=1:p, \quad \text{then} \quad c_i=0, \quad \forall\, i=1:p.
\]
\end{definition}

If $A$ is an $m\times n$ matrix, the column rank of $A$ is the largest
number of linearly independent columns of $A$, and the row rank of $A$ is
the largest number of linearly independent rows of $A$. However, the
column rank and the row rank of a matrix are always the same, not only for
square matrices, but also for rectangular matrices; see section 10.5 for an
elegant proof of this fact. Therefore, the following definition is
consistent:

\begin{definition}
The rank of a matrix $A$ is denoted by $\rank(A)$ and is equal to the
largest number of linearly independent columns\footnote{As explained
above, the rank of a matrix can also be defined as the largest number of
linearly independent rows of the matrix.} of $A$.
\end{definition}

\begin{definition}
The nullspace of an $m\times n$ matrix $A$ is
\[
\mathrm{Null}(A) \;=\; \{v\in\R^n \ \text{such that}\ Av=0\}.
\]
\end{definition}

The nullspace of the matrix $A$ is also called the kernel of $A$ and is
denoted by $\mathrm{Ker}(A)$.

\begin{definition}
The range of an $m\times n$ matrix $A$ contains all the vectors obtained
from matrix--vector multiplications involving $A$, i.e.,
\begin{equation}
\mathrm{Range}(A) \;=\; \{u\in\R^m \ \text{such that there exists}\ v\in\R^n
\ \text{with}\ u=Av\}.
\end{equation}
\end{definition}

\begin{definition}
Let $w_1,w_2,\ldots,w_p$ be vectors of the same size. The space generated
by $w_1, w_2, \ldots, w_p$ is made of all the linear combinations of the
vectors $w_1, w_2, \ldots, w_p$, i.e.,
\begin{equation}
<w_1,w_2,\ldots,w_p> \;=\; \left\{\sum_{i=1}^p c_iw_i, \quad \text{with}\
c_i\in\R,\ i=1:p\right\}.
\end{equation}
\end{definition}

\begin{lemma}
Let $A = \mathrm{col}\,(a_k)_{k=1:n}$ be an $m\times n$ matrix. The range
of the matrix $A$ is the space generated by the column vectors of $A$,
i.e.,
\begin{equation}
\mathrm{Range}(A) \;=\; <a_1,a_2,\ldots,a_n>.
\end{equation}
\end{lemma}

\begin{proof}
We prove (1.47) by double inclusion.

Let $u \in <a_1,a_2,\ldots,a_n>$. By definition (1.46), there exist
$c_i\in\R$, $i=1:n$, such that
\begin{equation}
u \;=\; \sum_{i=1}^n c_ia_i.
\end{equation}

Denote by $v=(c_i)_{i=1:n}$ the column vector with entries equal to the
coefficients from the linear combination (1.48). From (1.7), it follows
that
\begin{equation}
Av \;=\; \sum_{i=1}^n v_ia_i \;=\; \sum_{i=1}^n c_ia_i.
\end{equation}

From (1.48) and (1.49), we obtain that $u=Av$. Therefore, $u \in
\mathrm{Range}(A)$, and we conclude that
\begin{equation}
<a_1,a_2,\ldots,a_n> \;\subseteq\; \mathrm{Range}(A).
\end{equation}

Let $u \in \mathrm{Range}(A)$. From (1.45), it follows that there exists a
vector $v\in\R^n$ such that $u=Av$. Then, from (1.7) and (1.46), we obtain
that
\[
u \;=\; Av \;=\; \sum_{i=1}^n v_ia_i \;\in\; <a_1,a_2,\ldots,a_n>,
\]

and therefore
\begin{equation}
\mathrm{Range}(A) \;\subseteq\; <a_1,a_2,\ldots,a_n>.
\end{equation}

From (1.50) and (1.51), we conclude that $\mathrm{Range}(A) =
<a_1,a_2,\ldots,a_n>$.
\end{proof}

\begin{definition}
Let $V$ be a vector space in $\R^n$. The vectors $v_1,v_2,\ldots,v_p \in V$
form a basis for $V$ if they are linearly independent and if the space
generated by the vectors $v_1,v_2,\ldots,v_p$ is equal to $V$, i.e., if
$V = <v_1,v_2,\ldots,v_p>$.
\end{definition}

\begin{lemma}
Any two basis of a vector space in $\R^n$ are made of the same number of
vectors.
\end{lemma}

The proof of this result is of no further relevance in this book and can
be found, e.g., in Lax \cite{26}.

\begin{definition}
The dimension of a vector space $V$ from $\R^n$ is denoted by $\dim(V)$ and
is equal to the number of vectors from a basis of $V$.
\end{definition}

While a vector space has infinitely many different basis, Definition 1.11
is consistent, since, from Lemma 1.5, it follows that the numbers of
linearly independent vectors from any two different basis are equal.

\subsection{A one period market model}

In a one period market model for the evolution of the prices of financial
securities, portfolio returns can be analyzed using a linear algebra
framework. In this section, we present an example of this framework, which
forms the basis of the Arrow--Debreu model; see Chapter 3 for details on
the Arrow--Debreu model.

Consider a market with $m$ securities. Let $S_{1t_0}, S_{2t_0}, \ldots,
S_{mt_0}$ be the spot prices of the securities at time $t_0$, and denote by
$S_{t_0} = (S_{jt_0})_{j=1:m}$ the price vector of the securities at time
$t_0$. Note that $S_{t_0}$ is an $m\times 1$ column vector.

Assume that, at time $\tau > t_0$, there are $n$ possible states of the
market, denoted by $\omega^1, \omega^2, \ldots, \omega^n$. Let
$S_{j\tau}^{\omega^k}$ be the price at time $\tau$ of asset $j$ if state
$\omega^k$ occurs, for $1\le j\le m$ and $1\le k\le n$. Let

\[
S_\tau^{\omega^k} \;=\;
\begin{pmatrix} S_{1\tau}^{\omega^k} \\ S_{2\tau}^{\omega^k} \\ \vdots \\
S_{m\tau}^{\omega^k} \end{pmatrix}
\]

be the price vector of the $m$ assets if state $\omega^k$ occurs, for
$k=1:n$, and let

\[
S_{j\tau} \;=\; \left(S_{j\tau}^{\omega^1}\ \ S_{j\tau}^{\omega^2}\ \ \ldots\
\ S_{j\tau}^{\omega^n}\right)
\]

be the vector of all the possible prices of asset $j$ at time $\tau$, for
$j=1:m$. Note that $S_\tau^{\omega^k}$ is an $m\times 1$ column vector, and
$S_{j\tau}$ is an $1\times n$ row vector.

The payoff matrix $M_\tau$ is the $m\times n$ matrix made of all possible
asset prices at time $\tau$, with the $j$-th row of matrix $M_\tau$
corresponding to the prices of asset $j$, for $j=1:m$, and the $k$-th
column of matrix $M_\tau$ corresponding to the asset prices in state
$\omega^k$, for $k=1:n$, i.e.,

\begin{equation}
M_\tau \;=\; \mathrm{col}\,\left(S_\tau^{\omega^k}\right)_{k=1:n} \;=\;
\left(S_\tau^{\omega^1} \mid S_\tau^{\omega^2} \mid \ldots \mid
S_\tau^{\omega^n}\right);
\end{equation}

\begin{equation}
M_\tau \;=\; \mathrm{row}\,(S_{j\tau})_{j=1:m} \;=\;
\begin{pmatrix} S_{1\tau} \\ \text{--} \\ S_{2\tau} \\ \text{--} \\ \vdots \\
\text{--} \\ S_{m\tau} \end{pmatrix}.
\end{equation}

\medskip
\noindent\rule{\linewidth}{0.4pt}
\medskip

\noindent\textit{Example:}\ Consider two assets with spot prices \$30 and
\$50, respectively. Assume that, in three months, the first asset will be
worth either \$34 or \$24, and the second asset will be worth either \$56,
\$51, or \$46. The values of the three months at--the--money (ATM) European
call and put options\footnote{An at--the--money option has strike equal to
the spot price. A brief overview of European options can be found in
section 10.3; see also Neftci \cite{31} and Stefanica \cite{36}.}
with strike \$30 on the first asset are \$2.6 and \$2.3, respectively, and
the values of the three months ATM European call and put options with
strike \$50 on the second asset are \$2.7 and \$2.2, respectively.

Assume that the future value in three months of \$1 today is
\$1.01.\footnote{This corresponds, for continuous compounding, to a three
months risk free interest rate equal to $r=0.0398=3.98\%$, which can be
found by solving the equation $1.01 = 1\cdot\exp\left(r\cdot\frac{1}{4}\right)$
for $r$.}

A three months market model with seven securities and six states can be
constructed based on the information above, as follows:

\medskip
\noindent\underline{Securities:}
\begin{itemize}
\item cash;
\item first asset;
\item second asset;
\item three months ATM call with strike \$30 on the first asset;
\item three months ATM put with strike \$30 on the first asset;
\item three months ATM call with strike \$50 on the second asset;
\item three months ATM put with strike \$50 on the second asset.
\end{itemize}

\medskip
\noindent\underline{States of the market} in three months:
\begin{itemize}
\item first asset at \$34 and second asset at \$56 (state $\omega^1$);
\item first asset at \$34 and second asset at \$51 (state $\omega^2$);
\item first asset at \$34 and second asset at \$46 (state $\omega^3$);
\item first asset at \$24 and second asset at \$56 (state $\omega^4$);
\item first asset at \$24 and second asset at \$51 (state $\omega^5$);
\item first asset at \$24 and second asset at \$46 (state $\omega^6$).
\end{itemize}

This is a market model with $m=7$ securities and with $n=6$ states.

The price vector of the securities at time 0, corresponding to a \$1 cash
position and to positions equal to one unit of each of the other
securities, is

\begin{equation}
S_0 \;=\; \begin{pmatrix} 1 \\ 30 \\ 50 \\ 2.6 \\ 2.3 \\ 2.7 \\ 2.2 \end{pmatrix}.
\end{equation}

For $j=1:7$, let $S_{j,1/4}$ be the vector of the six possible prices of
asset $j$ in three months. The price vectors $S_{1,1/4}$ of cash,
$S_{2,1/4}$ of the first asset, and $S_{3,1/4}$ of the second asset,
respectively, are

\begin{align}
S_{1,1/4} &= (1.01\ \ 1.01\ \ 1.01\ \ 1.01\ \ 1.01\ \ 1.01); \\
S_{2,1/4} &= (34\ \ 34\ \ 34\ \ 24\ \ 24\ \ 24); \\
S_{3,1/4} &= (56\ \ 51\ \ 46\ \ 56\ \ 51\ \ 46).
\end{align}

To compute the price vectors $S_{4,1/4}, S_{5,1/4}, S_{6,1/4}, S_{7,1/4}$ of
the options, note that all the options have three months maturity, and
recall from (10.90) and (10.91) that the payoffs at maturity of call and
put options are

\begin{equation}
C(T) \;=\; \max(S(T)-K,0) \;=\; \begin{cases} S(T)-K, & \text{if}\ S(T) > K; \\
0, & \text{if}\ S(T) \le K; \end{cases}
\end{equation}

\begin{equation}
P(T) \;=\; \max(K-S(T),0) \;=\; \begin{cases} 0, & \text{if}\ S(T) \ge K; \\
K-S(T), & \text{if}\ S(T) < K. \end{cases}
\end{equation}

From (1.58), we obtain that the values of the three months ATM call with
strike \$30 on the first asset are given by

\[
C_1(1/4) \;=\; \max(S_1(1/4)-30,0),
\]

% UNCLEAR: page 30 says $S_1(1/4)$ "denotes the value of the second asset",
% but $S_1$ is the first asset throughout this example (strike \$30, values
% 34/24). Author's error; left as printed.
where $S_1(1/4)$ denotes the value of the second asset in three months,
and are as follows:

\begin{center}
\begin{tabular}{lll}
state $\omega^1$: & $S_1(1/4)=34$ and & $C_1(1/4)=\max(34-30,0)=4$; \\
state $\omega^2$: & $S_1(1/4)=34$ and & $C_1(1/4)=\max(34-30,0)=4$; \\
state $\omega^3$: & $S_1(1/4)=34$ and & $C_1(1/4)=\max(34-30,0)=4$; \\
state $\omega^4$: & $S_1(1/4)=24$ and & $C_1(1/4)=\max(24-30,0)=0$; \\
state $\omega^5$: & $S_1(1/4)=24$ and & $C_1(1/4)=\max(24-30,0)=0$; \\
state $\omega^6$: & $S_1(1/4)=24$ and & $C_1(1/4)=\max(24-30,0)=0$. \\
\end{tabular}
\end{center}

Thus, the vector $S_{4,1/4}$ of the six possible prices of the three months
ATM call with strike \$30 on the first asset is

\begin{equation}
S_{4,1/4} \;=\; (4\ \ 4\ \ 4\ \ 0\ \ 0\ \ 0).
\end{equation}

From (1.59), we obtain that the values of the three months ATM put with
strike \$30 on the first asset are given by

\[
P_1(1/4) \;=\; \max(30-S_1(1/4),0),
\]

and are as follows:

\begin{center}
\begin{tabular}{lll}
state $\omega^1$: & $S_1(1/4)=34$ and & $P_1(1/4)=\max(30-34,0)=0$; \\
state $\omega^2$: & $S_1(1/4)=34$ and & $P_1(1/4)=\max(30-34,0)=0$; \\
state $\omega^3$: & $S_1(1/4)=34$ and & $P_1(1/4)=\max(30-34,0)=0$; \\
state $\omega^4$: & $S_1(1/4)=24$ and & $P_1(1/4)=\max(30-24,0)=6$; \\
state $\omega^5$: & $S_1(1/4)=24$ and & $P_1(1/4)=\max(30-24,0)=6$; \\
state $\omega^6$: & $S_1(1/4)=24$ and & $P_1(1/4)=\max(30-24,0)=6$. \\
\end{tabular}
\end{center}

Thus, the vector $S_{5,1/4}$ of the six possible prices of the three months
ATM put with strike \$30 on the first asset is

\begin{equation}
S_{5,1/4} \;=\; (0\ \ 0\ \ 0\ \ 6\ \ 6\ \ 6).
\end{equation}

From (1.58), we obtain that the values of the three months ATM call with
strike \$50 on the second asset are given by

\[
C_2(1/4) \;=\; \max(S_2(1/4)-50,0),
\]

where $S_2(1/4)$ denotes the value of the second asset in three months, and
are as follows:

\begin{center}
\begin{tabular}{lll}
state $\omega^1$: & $S_2(1/4)=56$ and & $C_2(1/4)=\max(56-50,0)=6$; \\
state $\omega^2$: & $S_2(1/4)=51$ and & $C_2(1/4)=\max(51-50,0)=1$; \\
state $\omega^3$: & $S_2(1/4)=46$ and & $C_2(1/4)=\max(46-50,0)=0$; \\
state $\omega^4$: & $S_2(1/4)=56$ and & $C_2(1/4)=\max(56-50,0)=6$; \\
state $\omega^5$: & $S_2(1/4)=51$ and & $C_2(1/4)=\max(51-50,0)=1$; \\
state $\omega^6$: & $S_2(1/4)=46$ and & $C_2(1/4)=\max(46-50,0)=0$. \\
\end{tabular}
\end{center}

Thus, the vector $S_{6,1/4}$ of the six possible prices of the three months
ATM call with strike \$50 on the second asset is

\begin{equation}
S_{6,1/4} \;=\; (6\ \ 1\ \ 0\ \ 6\ \ 1\ \ 0).
\end{equation}

From (1.59), we obtain that the values of the three months ATM put with
strike \$50 on the second asset are given by

\[
P_2(1/4) \;=\; \max(50-S_2(1/4),0),
\]

and are as follows:

\begin{center}
\begin{tabular}{lll}
state $\omega^1$: & $S_2(1/4)=56$ and & $P_2(1/4)=\max(50-56,0)=0$; \\
state $\omega^2$: & $S_2(1/4)=51$ and & $P_2(1/4)=\max(50-51,0)=0$; \\
state $\omega^3$: & $S_2(1/4)=46$ and & $P_2(1/4)=\max(50-46,0)=4$; \\
state $\omega^4$: & $S_2(1/4)=56$ and & $P_2(1/4)=\max(50-56,0)=0$; \\
state $\omega^5$: & $S_2(1/4)=51$ and & $P_2(1/4)=\max(50-51,0)=0$; \\
state $\omega^6$: & $S_2(1/4)=46$ and & $P_2(1/4)=\max(50-46,0)=4$. \\
\end{tabular}
\end{center}

Thus, the vector $S_{7,1/4}$ of the six possible prices of the three months
ATM put with strike \$50 on the second asset is

\begin{equation}
S_{7,1/4} \;=\; (0\ \ 0\ \ 4\ \ 0\ \ 0\ \ 4).
\end{equation}

Summarizing, we conclude from (1.55--1.57) and (1.60--1.63) that

\begin{align*}
S_{1,1/4} &= (1.01\ \ 1.01\ \ 1.01\ \ 1.01\ \ 1.01\ \ 1.01); \\
S_{2,1/4} &= (34\ \ 34\ \ 34\ \ 24\ \ 24\ \ 24); \\
S_{3,1/4} &= (56\ \ 51\ \ 46\ \ 56\ \ 51\ \ 46); \\
S_{4,1/4} &= (4\ \ 4\ \ 4\ \ 0\ \ 0\ \ 0); \\
S_{5,1/4} &= (0\ \ 0\ \ 0\ \ 6\ \ 6\ \ 6); \\
S_{6,1/4} &= (6\ \ 1\ \ 0\ \ 6\ \ 1\ \ 0); \\
S_{7,1/4} &= (0\ \ 0\ \ 4\ \ 0\ \ 0\ \ 4).
\end{align*}

The payoff matrix $M_{1/4}$ is the following $7\times 6$ matrix:

\begin{equation}
M_{1/4} \;=\;
\begin{pmatrix} S_{1,1/4} \\ S_{2,1/4} \\ S_{3,1/4} \\ S_{4,1/4} \\ S_{5,1/4} \\ S_{6,1/4} \\ S_{7,1/4} \end{pmatrix}
\;=\;
\begin{pmatrix}
1.01 & 1.01 & 1.01 & 1.01 & 1.01 & 1.01 \\
34 & 34 & 34 & 24 & 24 & 24 \\
56 & 51 & 46 & 56 & 51 & 46 \\
4 & 4 & 4 & 0 & 0 & 0 \\
0 & 0 & 0 & 6 & 6 & 6 \\
6 & 1 & 0 & 6 & 1 & 0 \\
0 & 0 & 4 & 0 & 0 & 4
\end{pmatrix};
\end{equation}

see (1.53).

Let $S_{1/4}^{\omega^i}$ be the price vector of the seven securities in
three months if state $\omega^i$ occurs, for $i=1:6$. Recall from (1.52)
that

\begin{equation}
M_{1/4} \;=\; \left(S_{1/4}^{\omega^1}\mid S_{1/4}^{\omega^2}\mid
S_{1/4}^{\omega^3}\mid S_{1/4}^{\omega^4}\mid S_{1/4}^{\omega^5}\mid
S_{1/4}^{\omega^6}\right).
\end{equation}

From (1.64) and (1.65), we obtain that

\[
S_{1/4}^{\omega^1} = \begin{pmatrix}1.01\\34\\56\\4\\0\\6\\0\end{pmatrix}; \quad
S_{1/4}^{\omega^2} = \begin{pmatrix}1.01\\34\\51\\4\\0\\1\\0\end{pmatrix}; \quad
S_{1/4}^{\omega^3} = \begin{pmatrix}1.01\\34\\46\\4\\0\\0\\4\end{pmatrix};
\]

\[
S_{1/4}^{\omega^4} = \begin{pmatrix}1.01\\24\\56\\0\\6\\6\\0\end{pmatrix}; \quad
S_{1/4}^{\omega^5} = \begin{pmatrix}1.01\\24\\51\\0\\6\\1\\0\end{pmatrix}; \quad
S_{1/4}^{\omega^6} = \begin{pmatrix}1.01\\24\\46\\0\\6\\0\\4\end{pmatrix}.
\]

Consider a portfolio made of the $m$ securities and consisting of
$\Theta_j$ units of asset $j$ at time $t_0$, $j=1:m$. The value $V_{t_0}$
of the portfolio at time $t_0$ is

\[
V_{t_0} \;=\; \Theta_1S_{1t_0}+\Theta_2S_{2t_0}+\ldots+\Theta_mS_{mt_0},
\]

which can be written using the row vector--column vector multiplication
formula (1.5) as

\[
V_{t_0} \;=\; \Theta^tS_{t_0},
\]

where $\Theta=(\Theta_j)_{j=1:m}$ is the $m\times 1$ positions vector.

We assume that the asset positions in the portfolio remain unchanged until
time $\tau$. Let

\begin{equation}
V_\tau \;=\; \left(V_\tau^{\omega^1}\ \ V_\tau^{\omega^2}\ \ \ldots\ \
V_\tau^{\omega^n}\right)
\end{equation}

be the $1\times n$ row vector of the possible values of the portfolio at
time $\tau$, where $V_\tau^{\omega^k}$ is the value of the portfolio if
state $\omega^k$ occurs, for $k=1:n$. Then,

\begin{equation}
\begin{aligned}
V_\tau^{\omega^k} &= \Theta_1S_{1\tau}^{\omega^k}+\Theta_2S_{2\tau}^{\omega^k}
+\ldots+\Theta_mS_{m\tau}^{\omega^k} \\
&= \Theta^tS_\tau^{\omega^k},
\end{aligned}
\end{equation}

where (1.5) was used for deriving (1.67).

From (1.66) and (1.67), we obtain that

\begin{equation}
\begin{aligned}
V_\tau &= \left(V_\tau^{\omega^1}\ \ V_\tau^{\omega^2}\ \ \ldots\ \ V_\tau^{\omega^n}\right)
\;=\; \left(\Theta^tS_\tau^{\omega^1}\ \ \Theta^tS_\tau^{\omega^2}\ \ \ldots\ \ \Theta^tS_\tau^{\omega^n}\right) \\
&= \Theta^t\,\mathrm{col}\left(S_\tau^{\omega^k}\right)_{k=1:n} \\
&= \Theta^tM_\tau,
\end{aligned}
\end{equation}

since $M_\tau=\mathrm{col}\left(S_\tau^{\omega^k}\right)_{k=1:n}$; cf.\
(1.52).

\medskip
\noindent\rule{\linewidth}{0.4pt}
\medskip

\noindent\textit{Example (continued):}\ Consider the following portfolio in
the one period market model with seven securities and six states
introduced above:
\begin{itemize}
\item long \$10,000 cash;
\item long 200 units of the first asset;
\item short 100 units of the second asset;
\item short 1000 three months ATM calls on the first asset;
\item long 500 three months ATM calls on the second asset;
\item long 600 three months ATM puts on the second asset.
\end{itemize}

The positions vector of the portfolio is

\[
\Theta \;=\; \begin{pmatrix}10000\\200\\-100\\-1000\\0\\500\\600\end{pmatrix}.
\]

The value of the portfolio at time 0 is

\[
V_0 \;=\; \Theta^tS_0 \;=\; 11{,}070,
\]

where the price vector $S_0$ of the securities at time 0 is given by
(1.54).

Depending on the state of the market in three months, the possible values
of the portfolio are

\[
V_{1/4} \;=\; \Theta^tM_{1/4} \;=\; (10300\ \ 8300\ \ 10700\ \ 12300\ \ 10300\ \ 12700),
\]

where the payoff matrix $M_{1/4}$ is given by (1.64).

\medskip
\noindent\rule{\linewidth}{0.4pt}
\medskip

A derivative security is replicable in a market model with $m$ securities
and $n$ market states at time $\tau>t_0$ if there exists a portfolio made
of the $m$ securities which has the same values at time $\tau$ as the
derivative security in every state of the market at time $\tau$. Let
$s_\tau$ be the price vector at time $\tau$ of a replicable derivative
security, and let $\Theta=(\Theta_j)_{j=1:m}$ be the positions vector of
the replicating portfolio. If $V_\tau$ is the price vector of the
replicating portfolio at time $\tau$, then, from (1.68), we find that

\begin{equation}
s_\tau \;=\; V_\tau \;=\; \Theta^tM_\tau \;=\; \sum_{j=1}^m \Theta_jS_{j\tau},
\end{equation}

where for the last equality we used (1.9) and the fact that
$M_\tau=\mathrm{row}(S_{j\tau})_{j=1:m}$.

In other words, a derivative security is replicable if and only if the
price vector of the derivative security at time $\tau$ belongs to the
space generated by the price vectors $S_{1\tau}, S_{2\tau}, \ldots,
S_{m\tau}$ of the $m$ securities at time $\tau$; cf.\ Definition 1.9.

The $m$ securities are non-redundant if no security is replicable using
the other $m-1$ securities. Equivalently, the $m$ securities are
non-redundant if their price vectors at time $\tau$ are linearly
independent. Otherwise, the payoff of one of the securities would be
replicable by using the other $m-1$ securities, which would render the
security redundant. Thus, a one period market model has no redundant
securities if and only if the payoff matrix $M_\tau$ has rank equal to
$m$, i.e., $\rank(M_\tau)=m$.

\medskip
\noindent\rule{\linewidth}{0.4pt}
\medskip

\noindent\textit{Example (continued):}\ Recall that the price vectors in
three months of the seven securities from the market model considered
here are

\begin{align*}
S_{1,1/4} &= (1.01\ \ 1.01\ \ 1.01\ \ 1.01\ \ 1.01\ \ 1.01); \\
S_{2,1/4} &= (34\ \ 34\ \ 34\ \ 24\ \ 24\ \ 24); \\
S_{3,1/4} &= (56\ \ 51\ \ 46\ \ 56\ \ 51\ \ 46); \\
S_{4,1/4} &= (4\ \ 4\ \ 4\ \ 0\ \ 0\ \ 0); \\
S_{5,1/4} &= (0\ \ 0\ \ 0\ \ 6\ \ 6\ \ 6); \\
S_{6,1/4} &= (6\ \ 1\ \ 0\ \ 6\ \ 1\ \ 0); \\
S_{7,1/4} &= (0\ \ 0\ \ 4\ \ 0\ \ 0\ \ 4).
\end{align*}

Note that

\begin{equation}
S_{5,1/4}+S_{2,1/4}-S_{4,1/4} \;=\; (30\ \ 30\ \ 30\ \ 30\ \ 30\ \ 30)
\;=\; \frac{30}{1.01}S_{1,1/4}.
\end{equation}

Thus, the price vectors $S_{1,1/4}$, $S_{2,1/4}$, $S_{4,1/4}$, and
$S_{5,1/4}$ corresponding to cash, the first asset, the three months ATM
call and the three months ATM put on the first asset, respectively, are
not linearly independent.

We conclude that, e.g., the ATM put on the first asset can be replicated
using cash and positions on the first asset and ATM calls on the first
asset, since

\[
S_{5,1/4} \;=\; \frac{30}{1.01}S_{1,1/4} - S_{2,1/4} + S_{4,1/4},
\]

and therefore that the ATM put on the first asset is a redundant security.

Also, the price vectors $S_{1,1/4}$, $S_{3,1/4}$, $S_{6,1/4}$, and
$S_{7,1/4}$, corresponding to cash, the second asset, the three months ATM
call and the three months ATM put on the second asset, respectively, are
not linearly independent, since

\begin{equation}
S_{7,1/4}+S_{3,1/4}-S_{6,1/4} \;=\; (50\ \ 50\ \ 50\ \ 50\ \ 50\ \ 50)
\;=\; \frac{50}{1.01}S_{1,1/4}.
\end{equation}

We conclude that, e.g., the ATM put on the second asset can be replicated
using cash and positions on the second asset and ATM calls on the second
asset, since

\[
S_{7,1/4} \;=\; \frac{50}{1.01}S_{1,1/4} - S_{3,1/4} + S_{6,1/4},
\]

and therefore that the ATM put on the second asset is a redundant
security.

The redundancies above are due to the Put--Call parity,\footnote{For more
details on the Put--Call parity, see section 10.3 and section 1.9 from
Stefanica~\cite{36}.} which states that the values $C(t)$ and $P(t)$ of a
European call and of a European put option with the same strike $K$ and
maturity $T$ and on the same underlying asset with spot price $S(t)$
satisfy the following model independent relationship in arbitrage--free
markets:

\[
P(t)+S(t)-C(t) \;=\; Ke^{-r(T-t)}.
\]

Thus, at maturity,

\begin{equation}
P(T)+S(T)-C(T) \;=\; K.
\end{equation}

Then, (1.70) is a consequence of (1.72), since the vectors $S_{5,1/4}$,
$S_{2,1/4}$, $S_{4,1/4}$, and $S_{1,1/4}$ are the price vectors
corresponding to $P(T)$, $S(T)$, $C(T)$, and $K$, respectively, for the
three months ATM put and call on the first asset.

% UNCLEAR: page 35 ends this paragraph with "on the first asset", but
% $S_{3,1/4}$, $S_{6,1/4}$, $S_{7,1/4}$ are the second asset's vectors, so
% "second asset" was expected. Author's error; left as printed.
Similarly, (1.71) is a consequence of (1.72), since the vectors
$S_{7,1/4}$, $S_{3,1/4}$, $S_{6,1/4}$, and $S_{1,1/4}$ are the price
vectors corresponding to $P(T)$, $S(T)$, $C(T)$, and $K$, respectively, for
the three months ATM put and call on the first asset.
